\cleardoublepage
\addchap{Glossary}

This section is optional.
Ask your supervisor whether it is required for your thesis.

There are two ways to build a glossary:

\begin{itemize}
    \item If you just want to get it done, then use a \texttt{longtable} and fill your names in, see table below.
    \item If you want it fancy, then package \texttt{glossaries} (maybe \texttt{glossaries-extra}) may be your way to go.
    Be warned that although it automates symbol handling (e.g. sorting and referencing of symbols), it comes with some administrative overhead.
    You can find a discussion on different ways to achieve this \href{https://tex.stackexchange.com/a/366282}{on https://tex.stackexchange.com/a/366282}.
\item use the description-environment (see below)
\end{itemize}

\begin{center}
\begin{longtable}{@{}l p{10cm}@{}}
\toprule
Name & Description \\
\midrule
\endfirsthead
\multicolumn{2}{c}{\textit{Glossary -- continued}}\\
\toprule
Name & Description \\
\midrule
\endhead
\bottomrule \multicolumn{2}{r}{\textit{Continued on next page}} \\
\endfoot
\bottomrule
\endlastfoot

% start here with your Descriptions:
Actuator & Component for the conversion of an information carrying, low-energy control signal to a high energy signal capable of process intervention \\
Name2 & Description 2 \\
Name3 & Description 3 \\
\end{longtable}
\end{center}

\begin{description}
    \item [Actuator] Component for the conversion of an information carrying, low-energy control signal to a high energy signal capable of process intervention
    \item [Name2] Description 2
\end{description}
